\lab{2}{Поисковый робот}
\task
Разработать робота, принимающего единственный аргумент - путь к YAML-конфигу
с секциями \code{db} и \code{logic}. Сохранять нормализованный URL,
исходный HTML, источник и время загрузки Unix timestamp. Обеспечить
продолжение после остановки, периодическую проверку документов и обновление
содержимого только при изменении.

\method
Робот реализован на Python с \code{asyncio}, HTTPX и MongoDB.
Конфигурация задаёт соединение с базой, задержки, тайм-аут, период
повторной проверки и источники с начальными URL. Внутри источника
действуют разрешённые и исключённые префиксы пути. Навигационная страница
может служить отправной точкой, не попадая в итоговую коллекцию документов.

Нормализация URL приводит домен к нижнему регистру и IDNA, удаляет
стандартный порт и фрагмент, задаёт пустому пути значение \code{/}.
Параметры запроса сохраняются. Это устраняет часть повторов адресов,
но не решает общую задачу эквивалентности URL с разными query-параметрами.

Очередь хранится в MongoDB. Задания переходят между состояниями
\code{pending}, \code{processing}, \code{done}, \code{failed},
\code{blocked}. Атомарный захват задания предотвращает его одновременную
обработку. При новом запуске незавершённые задания возвращаются в очередь,
а завершённые перепроверяются после истечения заданного интервала.
Документ, найденные ссылки и завершение задания фиксируются транзакцией;
поэтому MongoDB запускается как одноузловой replica set.

Разные хосты обходятся параллельно, запросы к одному хосту - последовательно.
Задержка отсчитывается между началами запросов и учитывает ограничения
источника и \code{robots.txt}. Наибольшая задержка применяется ко всем
источникам одного хоста. Перенаправления обрабатываются явно, а новый адрес
повторно проверяется на допустимость обхода.

Условные HTTP-запросы используют ETag и Last-Modified. Ответ 304 не
перезаписывает содержимое. При ответе с телом сравнивается SHA-256:
одинаковый HTML повторно не сохраняется. Время последней проверки и время
изменения содержимого имеют разный смысл. Ошибка загрузки не должна удалять
уже сохранённую версию документа.

\journal
Логика разделена между конфигурацией, HTTP-клиентом, правилами источников,
хранилищем и движком обхода. После реализации последовательной загрузки
добавлены очередь в базе, восстановление и повторная проверка. Затем
введены ограничения \code{robots.txt}, общий регулятор задержек хоста
и конкурентный обход разных хостов. Сетевое ожидание вынесено из блокировки
операций БД, чтобы медленный сайт не останавливал остальные источники.

Аварийное завершение между отправкой HTTP-запроса и транзакцией может
привести к повторной загрузке страницы. Полностью исключить её нельзя без
транзакции с удалённым HTTP-сервером, но можно гарантировать отсутствие
потери задания и дублирования документа в базе.

\testing
\begin{longtable}{p{0.32\textwidth}p{0.60\textwidth}}
\toprule Сценарий & Ожидаемое поведение\\\midrule
Нормализация URL & Фрагмент удалён, query сохранён, стандартный порт убран.\\
Остановка и перезапуск & Задание processing снова доступно для обработки.\\
Повторный HTML, 304 & Тело документа не перезаписано.\\
Изменённый HTML & Новое содержимое и хеш сохранены.\\
Ошибки HTTP и тайм-аут & Старая версия сохранена, ошибка отмечена в очереди.\\
Запрет robots.txt & Документ не загружается в обход запрета.\\
Редирект за границы источника & Целевой адрес проверяется до загрузки.\\
Два источника одного хоста & Соблюдается общая задержка запросов.\\
Разные хосты & Сетевые ожидания могут перекрываться.\\\bottomrule
\end{longtable}
Эти группы представлены в \code{2_crawler/tests}: отдельные модули
проверяют очередь, хранение, повторный обход, robots, редиректы и задержки.
Автотесты запускаются командой \code{just test}. Они проверяют
правила и переходы состояния робота. Для проверки транзакций
с настоящей MongoDB нужен отдельный интеграционный запуск.

\results
Реализованы требуемый интерфейс конфигурации, возобновляемая очередь и
проверка изменений. Минимальный запуск:
\begin{lstlisting}
cd 2_crawler
docker compose up -d --wait
just run config.yaml
just test
\end{lstlisting}
Штатный пример задаёт задержку 2 секунды глобально и 3 секунды для Linux
Kernel, тайм-аут 30 секунд и повторную проверку через 86\,400 секунд.
Следовательно, задержка источника ограничивает частоту начала запросов
примерно одним запросом за три секунды без учёта более строгих правил сайта.
Это ограничение конфигурации, а не измеренная скорость выгрузки.

\conclusions
Сохранение очереди отдельно от процесса позволяет продолжать обход после
сбоя. Условные запросы сокращают передачу неизменившихся страниц,
а транзакции согласуют документ и очередь. Ограничения решения -
одно задание хоста за раз, отсутствие распределённой координации нескольких
роботов и необходимость MongoDB replica set. Большая часть итогового корпуса
получена импортом архивов, поэтому её объём нельзя считать результатом
HTTP-производительности робота.
